% Copyright (c) 2026 Christian Reitwiessner.
% Authors: Christian Reitwiessner
% Build: /tmp/andreev-pdf-tools/tectonic --keep-logs paper.tex
\documentclass[11pt]{article}
\usepackage[T1]{fontenc}
\usepackage{lmodern,amsmath,amssymb,amsthm,booktabs,microtype,enumitem}
\usepackage[margin=1in]{geometry}
\usepackage[colorlinks=true,allcolors=blue]{hyperref}
\newtheorem{theorem}{Theorem}[section]
\newtheorem{lemma}[theorem]{Lemma}
\newtheorem{proposition}[theorem]{Proposition}
\newtheorem{corollary}[theorem]{Corollary}
\newcommand{\bits}{\{0,1\}}
\newcommand{\ind}[1]{[\,#1\,]}
\newcommand{\THR}{\mathsf{THR}}
\newcommand{\smalltitle}[1]{\paragraph{#1}}
\setlength{\emergencystretch}{1em}
\setlist[itemize]{itemsep=1pt,topsep=3pt}
\setlist[enumerate]{itemsep=1pt,topsep=3pt}
\title{Exact and Approximate Threshold Circuits\\for Andreev's Function}
\author{Christian Reitwiessner}
\date{November 2026}
\hypersetup{pdfauthor={Christian Reitwiessner},
pdftitle={Exact and Approximate Threshold Circuits for Andreev's Function}}
\begin{document}
\maketitle
\begin{abstract}
\marginpar{\raggedright\scriptsize\textbf{TODO:} The abstract still needs considerable work.}
Andreev's function reads an input table at an address formed by parities
of disjoint input blocks. For the balanced $n$-input family, we construct
exact depth-two threshold circuits with $(5/16+o(1))n^2$ gates, improving
the $O(n^3/\log n)$ upper bound stated by Kane and Williams. Under
uniform inputs, explicit deterministic circuits achieve error at most
$\delta$ with $O(n^{3/2}\sqrt{\log(1/\delta)})$ gates, for
$0<\delta\le1/2$. At fixed $0<\delta<1/2$, the exponent matches
their lower bound up to logarithmic factors. Symmetric truncation of a
Fourier lookup formula makes the approximation errors odd in distinct
table coefficients; a table symmetry then makes these errors orthogonal.
Both constructions use bottom gates whose address weights have one sign
within each block. A restriction-counting argument gives an exact lower
bound of $\Omega(n^2/\log^3 n)$ for this class, with unrestricted real
table and output weights. We also give a size--advantage tradeoff and a
separate exponential bound for full-support, equal-magnitude address weights.
\end{abstract}

\section{Introduction}
For $d\in\mathbb Z_{\ge1}$, a threshold gate on $u\in\bits^d$ computes
$[\sum_{i=1}^d w_i u_i\ge\theta]$, where $w_i,\theta\in\mathbb R$
and $[P]\in\bits$ is the indicator of a statement $P$.
A depth-two threshold circuit consists of $G\in\mathbb Z_{\ge0}$ bottom
gates and one output threshold gate, which may also read the input bits
directly. All weights may be real. Size is $G+1$; wires are edges with
nonzero coefficients after coefficients on each edge are combined.
For a fixed circuit $C$, $C(u)$ denotes its output bit.
An affine map on a Boolean cube is the restriction of
$u\mapsto\beta_0+\sum_i\beta_i u_i$ with fixed real coefficients.

Fix integers $k,m\ge1$, and set $N=2^k$ and $n=N+km$.
Write $[j]=\{1,\ldots,j\}$ for $j\in\mathbb Z_{\ge0}$, with $[0]=\varnothing$.
The table space is $\mathcal X=\bits^{\bits^k}$, identified with $\bits^N$;
the address-input space is $\mathcal A=(\bits^m)^k$.
For $x=(x_c)_{c\in\bits^k}\in\mathcal X$ and
$a=(a_{i,j})_{i\in[k],j\in[m]}\in\mathcal A$, put
$|a_i|=\sum_{j=1}^m a_{i,j}$. Define the address map
$b:\mathcal A\to\bits^k$ and the ordinary Andreev function
$A_{k,m}:\mathcal X\times\mathcal A\to\bits$ by
$$
b_i(a)=|a_i|\bmod2,\qquad A_{k,m}(x,a)=x_{b(a)}.
$$
Table and address bits are separate inputs.
\emph{Floor balance} means $m=\lfloor N/k\rfloor$;
\emph{exact balance} means $km=N$ and is used only when $k\mid2^k$.
Balanced asymptotics let $k\to\infty$ through the respective admissible
integers, so $n\sim2N$. Unless stated otherwise, implicit constants are
absolute and uniform in all admissible parameters, including error and
advantage parameters. Logarithms are natural unless a base is displayed.
Empty sums are zero. Bold $\boldsymbol{x},\boldsymbol{a}$ denote independent
uniform random inputs in $\mathcal X,\mathcal A$; unbold $x,a$ are fixed
inputs. Uniform-input agreement of $C$ means
$\Pr[C(\boldsymbol{x},\boldsymbol{a})=A_{k,m}(\boldsymbol{x},\boldsymbol{a})]$;
error is its complement.

\paragraph{Results.}
Table~\ref{tab:bounds} summarizes the main comparisons.
A bottom gate is \emph{sign-coherent} if its weights within each address
block are all nonnegative or all nonpositive; signs may vary by block and
gate. A circuit is sign-coherent when all its bottom gates are.
Our upper bounds lie in this class. At balance its exact gate complexity
has exponent $2$, and its gate complexity at each fixed real error
$\delta\in(0,1/2)$ has exponent $3/2$, up to logarithms.

\begin{table}[t]
\centering
\small
\setlength{\tabcolsep}{4pt}
\renewcommand{\arraystretch}{1.25}
\begin{tabular}{@{}lll@{}}
\toprule
Setting and resource & Earlier bound~\cite{KW} & Bound here \\
\midrule
Exact, gate upper bound
 & $O(n^3/\log n)$ & $(5/16+o(1))n^2$ \\
Advantage $\varepsilon$, gates
 & $\Omega\!\left(\frac{\varepsilon^3n^{3/2}}{\log^3 n}\right)$
 & $O\!\left(\varepsilon^2n^{3/2}\sqrt{\log(1/\varepsilon)}\right)$ \\
Advantage $\varepsilon$, wires
 & $\Omega\!\left(\frac{\varepsilon^3n^{5/2}}{\log^{7/2}n}\right)$
 & $O\!\left(\varepsilon^2n^{5/2}\sqrt{\log(1/\varepsilon)}\right)$ \\
Exact, coherent lower bound
 & $\Omega(n^{3/2}/\log^3 n)$ & $\Omega(n^2/\log^3 n)$ \\
\bottomrule
\end{tabular}
\caption{Balanced ordinary Andreev lookup. Advantage $\varepsilon$ means
uniform-input agreement at least $1/2+\varepsilon$. The approximation rows
compare earlier lower bounds with our upper bounds, matching up to logarithms
at fixed advantage. Earlier bounds require $\varepsilon\gg\sqrt{\log n/n}$;
our displayed upper bounds hold for $N^{-1/2}\le\varepsilon\le1/8$.
Theorem~\ref{thm:avgupper} also handles larger advantage and gives the error
parameterization. All upper bounds here are sign-coherent. The last row
compares our restricted lower bound with the earlier unrestricted bound.}
\label{tab:bounds}
\end{table}

\paragraph{Relation to earlier work.}
Kane and Williams~\cite[Theorems~1.1--1.2]{KW} give the earlier bounds
in Table~\ref{tab:bounds}. Their Section~5.2 also gives a quadratic bound
in the different $\mathsf{MOD}_3\circ\mathsf{MOD}_2$ basis. Their threshold
construction expands lookup tests into parities and simulates them
separately; ours first groups by address character. Their Section~4
approximates parity by central-window staircases, and Section~6 notes
correlation with table majority. Our approximation adds odd-error
orthogonality and applies it to partial table votes.

Threshold staircases are classical; see Hajnal et al.~\cite[Proposition~2.1]{HMPST}
for symmetric functions and Hansen--Podolskii~\cite[Theorem~5]{HP} for a
related representation using bounded integer forms. Our sign-coherent
lower bound combines an information count under restrictions with the
single-crossing principle used in threshold lower bounds~\cite{IPS}.
It requires no weight bound. The polynomial overhead in simulations
with monotone bottom gates~\cite[Lemma~6]{HP} does not preserve our bounds.
The generalized Andreev function of Chen, Santhanam and
Srinivasan~\cite[Section~2.4]{CSS} uses an encoded table and an extractor
address, unlike the ordinary lookup function here. The more recent
$n^{5/2-\eta}$ gate lower bounds for functions in
$E^{\mathsf{NP}}$~\cite{CTW} also concern other functions; our exponent
gap is specific to ordinary Andreev lookup.

Appendix~\ref{sec:fullsign} gives a separate exponential lower bound
for full-support, equal-magnitude address weights. It uses the classical
parity-order argument of Minsky and Papert~\cite[Section~3.1]{MP}, which
forces a predicate to depend on every parity variable. It does not
improve the unrestricted lower bound. We make no exhaustive priority
claim for the elementary ingredients.

\section{Compiling functions of affine forms}\label{sec:identities}
\begin{lemma}[Affine-form compilation]\label{lem:ladder}
Let $d\in\mathbb Z_{\ge1}$, let $L:\bits^d\to\mathbb Z$ be affine,
and let $\ell,h\in\mathbb Z$ satisfy $\ell\le L(u)\le h$ for every
$u\in\bits^d$. For every function $g:\{\ell,\ldots,h\}\to\mathbb R$
and every $u\in\bits^d$,
$$
g(L(u))=g(\ell)+\sum_{\theta=\ell+1}^{h}
\bigl(g(\theta)-g(\theta-1)\bigr)[L(u)\ge\theta].
$$
More generally, fix $J\in\mathbb Z_{\ge0}$, an affine map
$B:\bits^d\to\mathbb R$, and, for each $j\in[J]$, an affine map $L_j:\bits^d\to\mathbb Z$,
integers $\ell_j,h_j$ with $\ell_j\le L_j(u)\le h_j$ for every $u\in\bits^d$,
and a function $g_j:\{\ell_j,\ldots,h_j\}\to\mathbb R$.
The function
$u\mapsto[B(u)+\sum_{j=1}^J g_j(L_j(u))\ge0]$
has a depth-two threshold circuit with direct input wires to the output
and at most
$\sum_{j=1}^J|\{\theta\in\mathbb Z:\ell_j<\theta\le h_j,\
g_j(\theta)\ne g_j(\theta-1)\}|$ bottom gates.
Identical comparison gates may be shared.
\end{lemma}
\begin{proof}
Fix $u\in\bits^d$. The active comparisons are exactly those with
$\theta\le L(u)$, so the sum telescopes.
Each comparison is a threshold gate on $u$; its output coefficient is
the fixed real number $g(\theta)-g(\theta-1)$.
The output threshold absorbs all constants, and direct input wires
supply the linear part of $B$.
\end{proof}

For each $S\subseteq[k]$, define affine maps
$X_S:\mathcal X\to\{-N,-N+2,\ldots,N\}$ and
$t_S:\mathcal A\to\{0,\ldots,m|S|\}$ by
$$
X_S(x)=\sum_{c\in\bits^k}(-1)^{S\cdot c}(1-2x_c),\qquad
t_S(a)=\sum_{i\in S}|a_i|,
$$
where $S\cdot c=\sum_{i\in S}c_i\in\mathbb Z$.
The map $X_S$ is fixed by $k,S$ and uses only table bits;
$t_S$ is fixed by $k,m,S$ and uses only address bits. In particular, $X_S(x)$ is an even integer for every $x$.
Character orthogonality gives, for every $(x,a)\in\mathcal X\times\mathcal A$,
\begin{equation}\label{eq:lookup}
\sum_{S\subseteq[k]}X_S(x)(-1)^{t_S(a)}=N(-1)^{x_{b(a)}}.
\end{equation}
Indeed, after exchanging sums, the coefficient of $(-1)^{x_c}$ is
$\prod_{i=1}^k(1+(-1)^{c_i+b_i(a)})$, equal to $N$ for $c=b(a)$ and
zero otherwise. For the matrix interpretation, define
$v(x)=((-1)^{x_c})_{c\in\bits^k}\in\mathbb R^N$ and
$X(x)=(X_S(x))_{S\subseteq[k]}\in\mathbb R^N$.
The fixed $N\times N$ matrix $H$ has rows indexed by $S\subseteq[k]$,
columns by $c\in\bits^k$, and entries $H_{S,c}=(-1)^{S\cdot c}$. It satisfies
$X(x)=Hv(x)$, $H^{\mathsf T}H=NI_N$, and
$v(x)=N^{-1}H^{\mathsf T}X(x)$.
Thus the numbers $X_S(x)/N$ are the coordinates of the signed truth
table $c\mapsto(-1)^{x_c}$ in the parity basis.
Equation~\eqref{eq:lookup} recovers its $b(a)$th entry by Fourier
inversion; see also \cite[Chapter~1]{OD}.

Define $\varphi:\mathbb Z\to\mathbb Z$ by $\varphi(u)=(-1)^u u$.

\begin{theorem}[A quadratic exact circuit]\label{thm:simple}
For every pair of integers $k,m\ge1$, with $N=2^k$ and $n=N+km$,
$A_{k,m}$ has an exact sign-coherent circuit with at most
$2N(N-1)+kmN$ bottom gates and one output gate.
All weights and thresholds can be integers of magnitude polynomial
in $n$, with an absolute degree and constant.
At floor or exact balance, the resulting size bound is $(3/4+o(1))n^2$.
\end{theorem}
\begin{proof}
For every even integer $v$ and integer $t$,
$$v(-1)^t=\varphi(v+t)-\varphi(t).$$
Apply this identity with $v=X_S(x)$ and $t=t_S(a)$.
For nonempty $S\subseteq[k]$, the affine maps
$(x,a)\mapsto X_S(x)+t_S(a)$ and $(x,a)\mapsto t_S(a)$ have integer
bounds $[-N,N+m|S|]$ and $[0,m|S|]$, of lengths
$2N+m|S|$ and $m|S|$. The empty-set term $X_\varnothing(x)$ is affine.
Lemma~\ref{lem:ladder} and $\sum_{S\subseteq[k]}|S|=kN/2$ give the
bottom-gate bound. Thresholding the negative of the left-hand side of
\eqref{eq:lookup} at zero computes $A_{k,m}$ on every input.
Bottom address weights lie in $\{0,1\}$, bottom input coefficients
have magnitude at most $2$, and bottom thresholds have magnitude $O(n)$.
The relevant finite differences of $\varphi$ have magnitude $O(n)$;
combining repeated gates preserves polynomial bounds.
\end{proof}

\section{Approximation from orthogonal odd errors}\label{sec:approx}
\begin{theorem}[Uniform-input approximation]\label{thm:avgupper}
There is an absolute constant $K>0$ such that, for every integer $k\ge1$
and $m=\lfloor2^k/k\rfloor$ (also covering exact balance), with
$N=2^k$ and $n=N+km$, the following hold:
\begin{enumerate}
\item For every real $0<\delta\le1/2$, an explicit deterministic
 sign-coherent depth-two circuit has error at most $\delta$ and
 at most $K n^{3/2}\sqrt{\log(1/\delta)}$ gates.
\item For every real $N^{-1/2}\le\varepsilon\le1/8$, an explicit
 deterministic sign-coherent depth-two circuit has agreement at least
 $1/2+\varepsilon$, at most
 $K\varepsilon^2n^{3/2}\sqrt{\log(1/\varepsilon)}$ gates, and at most
 $K\varepsilon^2n^{5/2}\sqrt{\log(1/\varepsilon)}$ wires.
\end{enumerate}
For every real $1/8<\varepsilon<1/2$, Part~1 applies with
$\delta=1/2-\varepsilon$.
\end{theorem}

\begin{lemma}[Odd-error orthogonality]\label{lem:odd}
Fix an integer $k\ge1$ and put $N=2^k$. For distinct $S,T\subseteq[k]$
and functions
$g,h:\mathbb Z\to\mathbb R$ with $g(-u)=-g(u)$ for every $u\in\mathbb Z$,
$$\mathbb E\bigl[g(X_S(\boldsymbol{x}))h(X_T(\boldsymbol{x}))\bigr]=0.$$
Here $\boldsymbol{x}$ is uniform in $\mathcal X$.
The assertion holds separately for each choice of $g,h$, including
choices indexed by a fixed address input $a\in\mathcal A$.
\end{lemma}
\begin{proof}
Choose $u\in\bits^k$ with $(S\mathbin\triangle T)\cdot u$ odd.
The map $T_u:\mathcal X\to\mathcal X$ defined by
$(-1)^{(T_u(x))_c}=(-1)^{T\cdot u}(-1)^{x_{c\oplus u}}$
is a bijection. It satisfies
$X_S(T_u(x))=-X_S(x)$ and $X_T(T_u(x))=X_T(x)$ for every $x$.
Thus $T_u(\boldsymbol{x})$ is uniform, and the expectation equals its
own negative. The finite table space guarantees that it exists.
\end{proof}
The Walsh coefficient values are not independent random variables;
the lemma uses their pairwise sign symmetry.

\begin{proof}[Proof of Theorem~\ref{thm:avgupper}]
\emph{Partial lookup votes.}
Fix admissible $k,m$, choose an integer $r\in\{0,\ldots,k\}$, and set
$q=2^r$ and $L=N/q\in\mathbb Z_{\ge1}$.
With $k,m,r$ fixed, define $P_*:\mathcal X\times\mathcal A\to\mathbb Z$ by
$$
P_*(x,a)=2\sum_{S\subseteq[r]}X_S(x)(-1)^{t_S(a)}
 =2q\sum_{\substack{c\in\bits^k\\c_{[r]}=b(a)_{[r]}}}(-1)^{x_c}.
$$
The second equality holds for every $x,a$ by character orthogonality;
$c_{[r]}$ denotes the first $r$ coordinates of $c$.
For each fixed $a$, the sum contains $L$ table entries, including $x_{b(a)}$.
Set the integer threshold $\theta_q=2q$ if $q<N$, and $\theta_q=0$ if
$q=N$. Define $V_r:\mathcal X\times\mathcal A\to\bits$ by
$V_r(x,a)=[-P_*(x,a)\ge\theta_q]$.
When $q=N$, $V_r=A_{k,m}$. Otherwise $L$ is even, and $V_r$ is majority
of the selected bits with ties outputting zero. Its agreement with
$A_{k,m}$, conditional on any fixed $a$, is $1/2+p(L)/2$, where
$p:2\mathbb Z_{\ge1}\to(0,1]$ is given by
$$
p(L)=2^{-L}\binom L{L/2}\ge(2L)^{-1/2}\qquad(L\in2\mathbb Z_{\ge1}).
$$
Indeed, conditioning on the other $L-1$ independent uniform table
signs, agreement exceeds $1/2$ only when their sum is $-1$; then it
is one. This event has probability $p(L)$.
The inequality follows by induction from $p(2)=1/2$ and
$p(L+2)/p(L)=(L+1)/(L+2)\ge\sqrt{L/(L+2)}$.

\emph{Symmetric truncation.}
For each nonempty $S\subseteq[r]$, define the affine map
$D_S:\mathcal A\to\mathbb Z$ and the constant $\eta_S\in\{-1,1\}$ by
$$D_S(a)=t_S(a)-\lfloor m|S|/2\rfloor,\qquad
\eta_S=(-1)^{\lfloor m|S|/2\rfloor}.$$
Both depend on $m,S$, and only $D_S(a)$ depends on the input.
For every $x,a$, evenness of $X_S(x)$ gives
$$
\eta_S\bigl(\varphi(X_S(x)+D_S(a))+\varphi(X_S(x)-D_S(a))\bigr)
 =2X_S(x)(-1)^{t_S(a)}.
$$
For an integer $R\ge1$, define $f_R:\mathbb Z\to\mathbb Z$ by
$f_R(u)=\varphi(u)[|u|\le R]$ and the truncated output expression
$P_R:\mathcal X\times\mathcal A\to\mathbb Z$ by
$$
P_R(x,a)=2X_\varnothing(x)+
\sum_{\varnothing\ne S\subseteq[r]}\eta_S
\bigl(f_R(X_S(x)+D_S(a))+f_R(X_S(x)-D_S(a))\bigr).
$$
For every integer $u$,
$$
f_R(u)=\sum_{\ell=-R}^{R+1}
\bigl(f_R(\ell)-f_R(\ell-1)\bigr)[u\ge\ell].
$$
Applying this identity to the two affine maps
$(x,a)\mapsto X_S(x)\pm D_S(a)$ gives at most $4(q-1)(R+1)$ bottom
gates for $P_R$. Together with the direct affine term and one output
gate, they form a circuit $C_{r,R}$ computing
$C_{r,R}(x,a)=[-P_R(x,a)\ge\theta_q]$.
Bottom address weights are all $+1$ or all $-1$ on the selected blocks
and zero elsewhere; hence $C_{r,R}$ is sign-coherent.

\emph{The mean-square estimate.}
For nonempty $S\subseteq[r]$, define $e_S:\mathcal X\times\mathcal A\to\mathbb Z$ by
\begin{align*}
e_S(x,a)={}&\eta_S\bigl(f_R(X_S(x)+D_S(a))+f_R(X_S(x)-D_S(a))\bigr)\\
&-2X_S(x)(-1)^{t_S(a)}.
\end{align*}
The maps depend on $k,m,S,R$ and sum to $P_R-P_*$ pointwise.
For each fixed $a$, the integer-valued map on $\mathbb Z$ given by
$$u\mapsto\eta_S\bigl((f_R-\varphi)(u+D_S(a))
                         +(f_R-\varphi)(u-D_S(a))\bigr)$$
is odd, since $f_R$ and $\varphi$ are odd; its value at $u=X_S(x)$
is $e_S(x,a)$. Apply Lemma~\ref{lem:odd} for each fixed $a$, then average
over $\boldsymbol a$, to obtain
\begin{equation}\label{eq:odd-l2}
\mathbb E\bigl[(P_R(\boldsymbol{x},\boldsymbol{a})-P_*(\boldsymbol{x},\boldsymbol{a}))^2\bigr]
 =\sum_{\varnothing\ne S\subseteq[r]}\mathbb E\bigl[e_S(\boldsymbol{x},\boldsymbol{a})^2\bigr].
\end{equation}

Fix such an $S$ and either choice of sign in the integer random variable
$U=X_S(\boldsymbol{x})\pm D_S(\boldsymbol{a})$.
Then $|\mathbb EU|\le1/2$. The centered random variable
$V=U-\mathbb EU$ sums $N$ independent uniform signs of magnitude $1$
and $m|S|$ independent uniform signs of magnitude $1/2$.
For the number $v_S=N+m|S|/4\ge2$, Hoeffding's bound gives
$\Pr(|V|>t)\le2e^{-t^2/(2v_S)}$ for every real $t\ge0$.
Since $R\ge1$, $|U|>R$ implies $|V|>R/2$, and $U^2\le2V^2+1/2$.
Integrating the tail bound yields
$$
\mathbb E\bigl[U^2[|U|>R]\bigr]
 \le(R^2+8v_S+1)e^{-R^2/(8v_S)}
 \le16v_Se^{-R^2/(16v_S)}.
$$
The last step uses $8v_S+1\le16v_S$ and $1+t\le e^t$ for $t\ge0$.
For every integer $u$, $|\varphi(u)-f_R(u)|=|u|[|u|>R]$.
Apply $(u+v)^2\le2u^2+2v^2$ to the two reflected errors and sum
\eqref{eq:odd-l2}. With $v=N+rm/4\ge v_S$, this gives
$$
\mathbb E\bigl[(P_R(\boldsymbol{x},\boldsymbol{a})-P_*(\boldsymbol{x},\boldsymbol{a}))^2\bigr]
 \le64(q-1)v e^{-R^2/(16v)}.
$$

\emph{From output expressions to agreement.}
If $q<N$ then $-P_*(x,a)\in4q\mathbb Z$, at distance
at least $2q$ from $\theta_q=2q$; if $q=N$, its value is $\pm2N$, at
distance $2q$ from $\theta_q=0$.
Thus $C_{r,R}(x,a)\ne V_r(x,a)$ implies
$|P_R(x,a)-P_*(x,a)|\ge2q$. Markov's inequality gives
$$\Pr[C_{r,R}(\boldsymbol{x},\boldsymbol{a})\ne V_r(\boldsymbol{x},\boldsymbol{a})]
 \le\frac{16v}{q}e^{-R^2/(16v)}.$$
For any real $0<\delta\le1/2$, the positive integer
$R=\lceil4\sqrt{v\log(16v/(q\delta))}\rceil$
depends only on $k,m,r,\delta$ and makes this probability at most $\delta$.
If this radius exceeds $N+km$, use $R=N+km$ instead: every attainable
$X_S(x)\pm D_S(a)$ then lies in $[-R,R]$, so $P_R=P_*$ everywhere.

For Part~1 choose $r=k$, so $q=N$ and $V_r=A_{k,m}$.
Balance gives $v\le5N/4$ and
$R=O(\sqrt{N\log(1/\delta)})$, uniformly for $0<\delta\le1/2$.
For Part~2 let $q$ be the least power of two at least
$32\varepsilon^2N$, choose $r=\log_2q$, and set $\delta=\varepsilon$.
The given range implies $32\le32\varepsilon^2N\le N/2$, so
$q<64\varepsilon^2N$ and $q\le N/2<N$.
The vote's advantage is at least $\sqrt{q/(8N)}\ge2\varepsilon$.
Losing at most $\varepsilon$ through truncation leaves advantage at least
$\varepsilon$. Since
$\log(16v/(q\varepsilon))=O(\log(1/\varepsilon))$ uniformly in this
range, the gate bound follows.
A circuit on $n$ inputs with $G\in\mathbb Z_{\ge0}$ bottom gates has at
most $nG+G+n$ wires, including direct input wires to the output.
The gate bound in Part~2 absorbs the additive terms in this wire count.
Taking one absolute $K$ large enough proves both parts.
\end{proof}

\section{Restrictions of sign-coherent circuits}\label{sec:coherent}
\begin{lemma}[Information preserved by a restriction]\label{lem:entropy}
Let $N,h\in\mathbb Z_{\ge1}$ and $G\in\mathbb Z_{\ge0}$.
Fix a depth-two threshold circuit $C$ on inputs
$(x,z)\in\bits^N\times\bits^h$, with $G$ bottom gates, arbitrary real weights, and direct input wires
to the output allowed. Let $\gamma_j:\bits^N\times\bits^h\to\bits$
be the output function of bottom gate $j\in[G]$.
Suppose the induced map
$$\bits^N\longrightarrow\bits^{\bits^h},\qquad
x\longmapsto\bigl(z\longmapsto C(x,z)\bigr)$$
is injective. Define $D:\bits^N\to\{0,\ldots,G\}$ by letting $D(x)$
count the bottom gates whose restricted output functions
$z\mapsto\gamma_j(x,z)$ are nonconstant.
For uniform $\boldsymbol{x}\in\bits^N$,
$$
\mathbb E[D(\boldsymbol{x})]\ge
\frac{\max\{0,N-h-1-\log_2(G+1)\}}{h+\log_2(G+1)}.
$$
\end{lemma}
\begin{proof}
For each fixed $x$, restricting a bottom gate shifts its threshold
without changing its weights on $z$. Varying this threshold yields at
most $2^h-1$ nonconstant Boolean functions on $\bits^h$.
For any integer $d\in\{0,\ldots,G\}$, specify the $d$ nonconstant gates
by their indices and one such function for each. Their output weights
and all direct $z$-weights are fixed by $C$.
Constant bottom outputs and direct table inputs only shift the output
threshold, leaving at most $2^h+1$ possible functions of $z$ for this
specification. Injectivity therefore gives
$$|\{x\in\bits^N:D(x)=d\}|\le\binom Gd(2^h-1)^d(2^h+1).$$
Writing $\mathsf H$ for base-two entropy, the chain rule and
$\binom Gd\le(G+1)^d$ yield
\begin{align*}
N=\mathsf H(\boldsymbol{x})
 &=\mathsf H(D(\boldsymbol{x}))+\mathsf H(\boldsymbol{x}\mid D(\boldsymbol{x}))\\
 &\le\log_2(G+1)+\mathbb E[D(\boldsymbol{x})]
       \bigl(\log_2(G+1)+h\bigr)+h+1.
\end{align*}
Rearrange and use $D(x)\ge0$ for every $x$.
\end{proof}

\begin{theorem}[Exact sign-coherent lower bound]\label{thm:coherent}
For every pair of integers $k,m\ge1$, with $N=2^k$ and $n=N+km$,
every exact sign-coherent circuit for $A_{k,m}$ with
$G\in\mathbb Z_{\ge0}$ bottom gates satisfies
$$
G\ge\frac{D_0}{1-(1-1/m)^k}\ge\frac m kD_0,
\qquad
D_0=\frac{\max\{0,N-k-1-\log_2(G+1)\}}{k+\log_2(G+1)}.
$$
Here $D_0$ is a nonnegative real number depending only on $k,G$.
At floor or exact balance, $G=\Omega(n^2/\log^3 n)$ as $k\to\infty$.
\end{theorem}
\begin{proof}
\emph{Each restriction reveals the table.}
Fix such a circuit $C$, and define
$\iota:\{0,\ldots,m\}^k\to\mathcal A$ by letting $\iota(s)$ have ones
in exactly the first $s_i$ positions of block $i$.
For every index $q\in\{0,\ldots,m-1\}^k$, the map
$z\mapsto\iota(q+z)$, $z\in\bits^k$, leaves one live bit per block, and
$$C(x,\iota(q+z))=A_{k,m}(x,\iota(q+z))
  =x_{(q\bmod2)\oplus z}\qquad(x\in\mathcal X,\ z\in\bits^k).$$
Here $q\bmod2$ is taken coordinatewise.
The map $x\mapsto(z\mapsto C(x,\iota(q+z)))$ is injective because its
values read every table entry. A bottom gate \emph{splits} this cube at
$x$ if its output is nonconstant as a function of $z$.
Lemma~\ref{lem:entropy}, with $h=k$, gives at least $D_0$ splitting gates
in expectation over $\boldsymbol{x}$ for each fixed $q$.

\emph{A gate splits few cubes.}
Fix $x\in\mathcal X$ and a bottom-gate index $j\in[G]$.
Let $F_{j,x}:\{0,\ldots,m\}^k\to\mathbb R$ be its weighted input sum
on $(x,\iota(s))$, and let $\theta_j\in\mathbb R$ be its threshold.
Sign coherence supplies a vector $\epsilon\in\{-1,1\}^k$ such that
$F_{j,x}$ is nondecreasing in direction $\epsilon_i$ in coordinate $i$.
For any fixed cube index $q$, its minimum and maximum occur at
$$q^-=q+(\mathbf1-\epsilon)/2,\qquad
q^+=q+(\mathbf1+\epsilon)/2,$$
where $\mathbf1$ is the all-one vector in $\mathbb Z^k$.
Thus the gate splits this cube exactly when
$F_{j,x}(q^-)<\theta_j\le F_{j,x}(q^+)$.
This argument permits unequal magnitudes and zero weights.

Partition the index grid into maximal lines with consecutive indices
differing by $\epsilon$. The maximum vertex of one cube is the minimum
vertex of the next, so the intervals
$(F_{j,x}(q^-),F_{j,x}(q^+)]$ along a line are disjoint.
The gate splits at most one cube per line.
A line starts exactly where $q-\epsilon$ is outside the grid; precisely
$(m-1)^k$ indices have a predecessor. Hence there are
$m^k-(m-1)^k$ lines, an upper bound on the number of cubes split by
this fixed gate at this fixed $x$.

\emph{Combining the counts.}
Sum the expected incidences of gates and split cubes over all $m^k$
cube indices. The pointwise gate bound gives
$$m^kD_0\le G\bigl(m^k-(m-1)^k\bigr).$$
Divide by the positive denominator and use
$1-(1-1/m)^k\le k/m$ to obtain the two inequalities.
For the asymptotic claim, $G\ge N^2$ already suffices. Otherwise,
integrality gives $G+1\le N^2$, and for all sufficiently large $k$,
$D_0\ge(N-3k-1)/(3k)$. At either balance,
$m=\Theta(N/k)$ and $n=\Theta(N)$, yielding the stated bound with
an absolute implicit constant.
\end{proof}

\paragraph{Why exactness matters.}
These restrictions range over block counts, not uniform raw addresses,
so a uniform-input error bound need not survive them. This allows the
sign-coherent approximation circuits of Section~\ref{sec:approx}.
For any fixed real $\eta>0$, an exact balanced circuit family of size
$O(n^{2-\eta})$ must, for all sufficiently large admissible $k$, use
mixed address-weight signs within a block of some bottom gate.

\section{Further questions}
For ordinary Andreev lookup, the unrestricted exact gate bounds still
have exponents $3/2$ and $2$. Can mixed signs within address blocks
yield $O(n^{2-\eta})$ gates for some fixed real $\eta>0$? Can the exact wire bound
of $O(n^3)$ improve toward the known lower exponent $5/2$? Gate upper
bounds imply wire upper bounds with a factor of $O(n)$; wire lower
bounds require information about which inputs each gate uses.

At any fixed real error $\delta\in(0,1/2)$, the gate exponent is determined up to
logarithmic factors, but the error dependence remains open. Our bound
becomes quadratic when $\log(1/\delta)=\Theta(n)$. For small advantage,
the upper bound's $\varepsilon^2$ dependence differs from the lower
bound's $\varepsilon^3$, where that lower bound applies. Closing either
gap would clarify the transition from weak prediction to exact computation.
Do structured or biased table distributions retain enough of the sign
symmetry in Lemma~\ref{lem:odd} for the same approximation argument?

\paragraph{Verification and assistance.}
The research and exposition were developed with AI assistance.
Finite checks accompany the source: \texttt{verify\_results.py} checks
the refined construction; \texttt{approximation\_checks.py} checks
truncation and orthogonality; and \texttt{restriction\_checks.py} checks
the restrictions. They supplement the proofs. An earlier quadratic
construction is formalized in Lean in
\href{https://github.com/cslib-community/cslib-frontier/pull/2813}{cslib-frontier};
no formalization is claimed for the other results here.

\appendix
\section{The sharper exact construction}\label{sec:exact}
Auxiliary notation is local to this appendix.

\begin{theorem}\label{thm:exact}
For every pair of integers $k,m\ge1$, put $N=2^k$, $n=N+km$, and
$\alpha=\lfloor(m-1)/2\rfloor\in\mathbb Z_{\ge0}$.
There is an exact sign-coherent depth-two threshold circuit for
$A_{k,m}$ with at most
\begin{equation}\label{eq:bestcount}
\frac{N(N-1)}2+\frac{kN(m+\alpha)}2
 +2(N-1)+2\alpha k+3km+\left\lceil\frac{km}{2}\right\rceil
\end{equation}
bottom gates and one output gate. All weights and thresholds are integers
of magnitude polynomial in $n$, with an absolute degree and constant.
At exact balance the resulting upper bound on total size is
$$\frac5{16}n^2-\frac n4\log_2(n/2)+O(n).$$
At floor balance it is $(5/16+o(1))n^2$ as $k\to\infty$.
\end{theorem}

\begin{proof}
\emph{Affine maps and the staircase.}
Fix $k,m$ as in the theorem. Let $0\in\bits^k$ be the zero address and
$e_i\in\bits^k$ the address with only coordinate $i\in[k]$ equal to one.
Define integer-valued affine maps $s_i,\sigma$ on $\mathcal A$ and
$Z,E$ on $\mathcal X$ by
$$s_i(a)=|a_i|,\quad \sigma(a)=\sum_{i=1}^k s_i(a),\quad
Z(x)=\sum_{c\ne0}x_c,\quad E(x)=\sum_{i=1}^k x_{e_i}.$$
Their respective ranges lie in $[0,m]$, $[0,km]$, $[0,N-1]$, and $[0,k]$.
Set the integer constants $R=2(N-1)+2\alpha k+km$ and $C_0=N(R+4)$.
Define affine maps $w:\mathcal X\times\mathcal A\to\{0,\ldots,R\}$ and,
for each nonempty $S\subseteq[k]$,
$W_S:\mathcal X\to\{0,\ldots,N/2+\alpha|S|\}$ by
$$w(x,a)=2Z(x)+2\alpha E(x)+\sigma(a),\qquad
W_S(x)=\sum_{\substack{c\in\bits^k\\S\cdot c\text{ odd}}}x_c
       +\alpha\sum_{i\in S}x_{e_i}.$$
Finally, define $\kappa_S=-4(2\alpha+1)(k-2|S|)\in\mathbb Z$ and
$r_S:\mathbb Z\to\mathbb Z$ by
$r_S(u)=-8\lceil u/2\rceil^2+\kappa_S\lceil u/2\rceil$.
All constants and maps are fixed by $k,m$ and their displayed indices;
$W_S,Z,E$ depend on table inputs, $s_i,\sigma,t_S$ on address inputs,
and $w$ on both. For every pair of integers $W,t$,
$$4W(-1)^t=4W-8\lceil(2W+t)/2\rceil^2+8W^2+8Wt+8\lceil t/2\rceil^2.$$
The function $r_S$ changes only at odd integer levels, saving thresholds;
$\kappa_S$ corrects singleton addresses using those same thresholds.

\emph{The output expression and its sign.}
Define $P:\mathcal X\times\mathcal A\to\mathbb Z$ by
\begin{align}\label{eq:bestscore}
P(x,a)={}&2N+\sum_{\varnothing\ne S\subseteq[k]}r_S(2W_S(x)+t_S(a))
 +\frac N2\left(w(x,a)^2+\sum_{i=1}^k s_i(a)^2\right)\notag\\
 &+4NZ(x)-2N(\alpha^2+1)E(x)-2\alpha N\sigma(a)\notag\\
 &+2N(1+\alpha)\sum_{i=1}^k x_{e_i}s_i(a)-C_0x_0\ind{b(a)=0}.
\end{align}
To verify its sign, fix arbitrary inputs $x\in\mathcal X,a\in\mathcal A$.
Only in the following algebra, use the integer abbreviations
$$u_i=s_i(a),\quad u=\sigma(a),\quad z=Z(x),\quad e=E(x),\quad
\omega=w(x,a),\quad \omega_S=W_S(x),\quad \tau_S=t_S(a),$$
and the address $\beta=b(a)\in\bits^k$.
These are values at the fixed inputs, not maps. In the next identities,
$\sum_S$ ranges over nonempty $S\subseteq[k]$ and $\sum_i$ over $i\in[k]$.
Expanding the staircase identity uses
\begin{align*}
\sum_S4\omega_S(-1)^{\tau_S}
 &=2N\ind{\beta=0}z-2N\ind{\beta\ne0}x_\beta
   +2\alpha N\sum_i x_{e_i}\bigl(\ind{\beta=0}-\ind{\beta=e_i}\bigr),\\
8\sum_S(\omega_S^2+\omega_S\tau_S)
 &=\frac N2(\omega^2-u^2)+2Nz+(4\alpha+2\alpha^2)Ne
   +2N(1+\alpha)\sum_i x_{e_i}u_i,\\
8\sum_S\lceil\tau_S/2\rceil^2
 &=\frac N2\left(u^2+\sum_i u_i^2\right)
   +N(u+1)\ind{\beta\ne0}+N\sum_i u_i\ind{\beta=e_i}.
\end{align*}
The first identity follows from character orthogonality.
For any distinct nonzero addresses $c,d\in\bits^k$,
$$\sum_S[S\cdot c\text{ odd}]=N/2,\qquad
\sum_S[S\cdot c\text{ odd}][S\cdot d\text{ odd}]=N/4;$$
when $c=d\ne0$, the second count is $N/2$.
These counts expand the quadratic terms.
For the last identity use $\lceil t/2\rceil=(t+[t\text{ odd}])/2$
for integers $t$, followed by the same orthogonality.
The shifts sum using $k-2|S|=\sum_i(-1)^{[i\in S]}$:
\begin{gather*}
\sum_S\kappa_S(\lceil\tau_S/2\rceil-\tau_S/2)
   =N(2\alpha+1)\sum_i\ind{\beta=e_i},\\
\sum_S\kappa_S\omega_S=2N(2\alpha+1)(1+\alpha)e,\qquad
\frac12\sum_S\kappa_S\tau_S=N(2\alpha+1)u.
\end{gather*}
Substituting in \eqref{eq:bestscore} and restoring function arguments gives,
for every $x,a$,
\begin{equation}\label{eq:bestclosed}
P(x,a)=\begin{cases}
N(w(x,a)+2)-C_0x_0,&b(a)=0,\\
2N(1+\alpha)(1-x_{e_i})-Ns_i(a),&b(a)=e_i\ (i\in[k]),\\
N(1-2x_{b(a)}),&\text{otherwise}.
\end{cases}
\end{equation}
If $b(a)=0$, use $0\le w(x,a)\le R$ and $C_0=N(R+4)$.
If $b(a)=e_i$, the odd integer $s_i(a)$ is between $1$ and $2\alpha+1$.
In all cases $P(x,a)\ge N$ when $x_{b(a)}=0$, and $P(x,a)\le-N$ when
$x_{b(a)}=1$. Hence $(x,a)\mapsto[-P(x,a)\ge0]$ computes $A_{k,m}$.

\emph{Realizing and counting the gates.}
We now use maps on the whole input cube, without the pointwise abbreviations.
For each nonempty $S\subseteq[k]$ the term
$r_S(2W_S(x)+t_S(a))$ uses the gates
$$\ind{2W_S(x)+t_S(a)\ge\theta},\qquad
\theta\in\mathbb Z\text{ odd},\quad
1\le\theta\le N+2\alpha|S|+m|S|,$$
with output coefficient $-8\theta+\kappa_S$ in the expansion of $P$.
Thus the shifts $\kappa_S$ add no gates.
Lemma~\ref{lem:ladder} supplies $R+km$ gates for
$w(x,a)^2$ and $\sum_{i=1}^k s_i(a)^2$.
With the integer $B=n+2>km$, each product map
$(x,a)\mapsto x_{e_i}s_i(a)$ is a sum of the $m$ gates
$[s_i(a)+Bx_{e_i}\ge B+\ell]$, $\ell\in\{1,\ldots,m\}$.

For the zero-address correction, the identity
$$x_0\ind{b(a)=0}=x_0-x_0\sigma(a)+
 \frac4N\sum_{\varnothing\ne S\subseteq[k]}x_0\lfloor t_S(a)/2\rfloor$$
holds for every $x,a$. It follows from
$\sum_{S\subseteq[k]}[t_S(a)\text{ odd}]=(N/2)[b(a)\ne0]$
and $\sum_{S\subseteq[k]}t_S(a)=N\sigma(a)/2$.
The product $x_0\sigma(a)$ uses $km$ gates
$[\sigma(a)+Bx_0\ge B+\ell]$, $\ell\in\{1,\ldots,km\}$.
For each nonempty $S$, the product $x_0\lfloor t_S(a)/2\rfloor$ uses
$[t_S(a)+Bx_0\ge B+2j]$,
$j\in\{1,\ldots,\lfloor m|S|/2\rfloor\}$.
The even-level $\sigma$ gates coincide with the latter gates for $S=[k]$,
saving $\lfloor km/2\rfloor$ gates.
Every remaining term of $P$ is affine and can be supplied directly to
the output. All these bottom gates have nonnegative address weights.

Since $\sum_{\varnothing\ne S\subseteq[k]}|S|=kN/2$, the character
staircases and the floor-product terms together use at most
$$\sum_{\varnothing\ne S\subseteq[k]}\left(
 \left\lceil\frac{N+2\alpha|S|+m|S|}{2}\right\rceil
 +\left\lfloor\frac{m|S|}{2}\right\rfloor\right)
 =\frac{N(N-1)}2+\frac{\alpha kN}2+\frac{kmN}2.$$
Adding the other gates and subtracting the shared ones gives
\eqref{eq:bestcount}; zero coefficients or further coincidences can
reduce the actual size.
At exact balance $m$ is even and $\alpha=m/2-1$, giving the claimed
upper-bound expansion after adding the output gate.
At floor balance, rounding changes the leading expression by
$O(Nk)=o(n^2)$.
Bottom coefficients and thresholds have magnitude $O(n)$; output
coefficients, including the output threshold, have magnitude $O(n^3)$.
Combining further repeated gates still leaves polynomial bounds.
All coefficients are integers, since $N$ is even and $C_0/N$ is an integer.
\end{proof}

\section{A full-support equal-magnitude obstruction}\label{sec:fullsign}
The following classical parity-order obstruction~\cite[Section~3.1]{MP}
yields the full-support, equal-magnitude lower bound.

\begin{lemma}[The missing-variable parity obstruction]\label{lem:paritydual}
Let $d\in\mathbb Z_{\ge2}$ and $J\in\mathbb Z_{\ge0}$, and define
$f:\{-1,1\}^d\to\{-1,1\}$ by $f(z)=\prod_{i=1}^d z_i$.
Fix predicates $p_j:\{-1,1\}^d\to\bits$ and coefficients
$\lambda_j\in\mathbb R$ for $j\in[J]$, and an affine map
$B:\{-1,1\}^d\to\mathbb R$,
$B(z)=\beta_0+\sum_{i=1}^d\beta_i z_i$, with $\beta_0,\ldots,\beta_d\in\mathbb R$.
Define $F:\{-1,1\}^d\to\mathbb R$ by
$F(z)=B(z)+\sum_{j=1}^J\lambda_jp_j(z)$.
If $f(z)F(z)>0$ for every $z\in\{-1,1\}^d$, then some $j\in[J]$
has $\lambda_j\ne0$ and $p_j$ depends on every coordinate: for each
$i\in[d]$ there exist inputs differing only in coordinate $i$ on which
$p_j$ has different values.
\end{lemma}
\begin{proof}
For uniform $\boldsymbol z\in\{-1,1\}^d$, a predicate independent of
coordinate $i$ has
$\mathbb E[f(\boldsymbol z)p_j(\boldsymbol z)]=0$, by summing over that
coordinate. Also $\mathbb E[f(\boldsymbol z)B(\boldsymbol z)]=0$
since $d\ge2$. If each predicate with nonzero coefficient omitted a
coordinate, linearity would give $\mathbb E[f(\boldsymbol z)F(\boldsymbol z)]=0$.
The assumed strict product inequality instead gives
$\mathbb E[f(\boldsymbol z)F(\boldsymbol z)]=\mathbb E|F(\boldsymbol z)|>0$.
\end{proof}

\begin{proposition}\label{thm:fullsign}
Let $k,m\in\mathbb Z$ satisfy $k\ge2$, $m\ge1$, and
$d=m\lfloor k/2\rfloor\ge2$, and put $N=2^k$, $n=N+km$.
Fix an exact depth-two threshold circuit for $A_{k,m}$ with
$G\in\mathbb Z_{\ge0}$ bottom gates.
For each $j\in[G]$, suppose its address-weight vector
$v_j\in\mathbb R^{km}$ is zero, or there exist a gate-specific real number
$\rho_j>0$ and vector $\sigma_j\in\{-1,1\}^{km}$ such that
$v_j=\rho_j\sigma_j$.
Then $G\ge2^d$. Table weights, thresholds, output weights, and direct
input wires to the output are unrestricted; all weights may be real.
At floor or exact balance this gives $G\ge2^{\Omega(n)}$ as $k\to\infty$.
\end{proposition}
\begin{proof}
Set $r=\lfloor k/2\rfloor$ and fix the table to
$x^*\in\mathcal X$ given by $x^*_c=(\sum_{i=1}^r c_i)\bmod2$ for every
$c\in\bits^k$. If $k$ is odd, fix one unselected block to zero.
There remain $d=rm$ selected bits and $d$ ignored bits.
Pair them, and use their sign coordinates $2a_{i,j}-1$ to obtain
$z,w\in\{-1,1\}^d$. The target bit is the parity of the selected bits,
so twice its value minus one is $(-1)^{d+1}f(z)$, where $f(z)=\prod_{i=1}^d z_i$.

After these fixings, write the circuit output as $[F_0(z,w)\ge0]$,
where $F_0:\{-1,1\}^{2d}\to\mathbb R$ is an affine function of $(z,w)$ plus a fixed real linear combination
of the bottom predicates,
with the output threshold subtracted.
On inputs of target value zero, $F_0(z,w)<0$. Finiteness allows a constant
$\gamma>0$ smaller than all the corresponding $|F_0(z,w)|$.
The map $F:\{-1,1\}^{2d}\to\mathbb R$ defined by
$F(z,w)=(-1)^{d+1}(F_0(z,w)+\gamma)$ satisfies
$f(z)F(z,w)>0$ for every $z,w$.

For each fixed $t\in\{-1,1\}^d$, impose $w_i=t_i z_i$ for all $i\in[d]$.
The direct output terms remain affine in $z$.
For a bottom gate $j$ with nonzero address weights, write
$s_{j,i},u_{j,i}\in\{-1,1\}$ for the two entries of $\sigma_j$
on pair $i$. Its restricted weighted sum has coefficient
$(\rho_j/2)(s_{j,i}+t_i u_{j,i})$ on $z_i$.
Dependence on every $z_i$ requires every such coefficient to be nonzero,
hence $t_i=s_{j,i}u_{j,i}$ for all $i$. This identifies at most one $t$
for that gate; nonzero coefficients alone need not ensure dependence.
A gate with zero address weights is constant after fixing the table.
Apply Lemma~\ref{lem:paritydual} to $z\mapsto F(z,(t_i z_i)_{i=1}^d)$:
each of the $2^d$ restrictions requires a bottom predicate depending on
all $d$ coordinates, and each original gate can serve at most one.
Thus $G\ge2^d$. At balance $d=\Theta(n)$, giving the asymptotic claim.
\end{proof}

Both address-weight qualifications are necessary.
Without full support, Theorem~\ref{thm:simple} gives quadratic size with
bottom address weights in $\{0,1\}$.
Without equal magnitudes, for every $S\subseteq[k]$ define the affine map
$L_S:\mathcal A\to\{0,\ldots,2km\}$ by
$$L_S(a)=\sum_{i\in S}\sum_{j=1}^m a_{i,j}
          +2\sum_{i\notin S}\sum_{j=1}^m a_{i,j}.$$
It has full address support and $L_S(a)\equiv t_S(a)\pmod2$ for every $a$.
Thus, for every $x,a$,
$$X_S(x)(-1)^{t_S(a)}
 =\varphi(X_S(x)+L_S(a))-\varphi(L_S(a)).$$
The affine arguments have integer interval lengths at most $2N+2km$
and $2km$. Applying Lemma~\ref{lem:ladder} to \eqref{eq:lookup} therefore
gives $O(Nn)=O(n^2)$ gates for every $k,m\ge1$, each bottom gate having
full address support with weights in $\{1,2\}$.
The exponential lower bound does not extend to full-support address
weights of general bounded integer magnitudes.

\begin{thebibliography}{9}\small
\bibitem{OD} R.~O'Donnell. \emph{Analysis of Boolean Functions}.
Cambridge University Press, 2014.
\href{https://www.cs.cmu.edu/~odonnell/papers/Analysis-of-Boolean-Functions-by-Ryan-ODonnell.pdf}{Author-hosted edition}.
\bibitem{CTW} L.~Chen, A.~Tal, and Y.~Wang.
Super-quadratic Lower Bounds for Depth-2 Linear Threshold Circuits.
\href{https://eccc.weizmann.ac.il/report/2026/039/}{ECCC TR26-039}, 2026.
\bibitem{KW} D.~M. Kane and R.~Williams.
Super-Linear Gate and Super-Quadratic Wire Lower Bounds for Depth-Two and
Depth-Three Threshold Circuits. STOC, 2016.
\href{https://arxiv.org/abs/1511.07860}{arXiv version: 1511.07860}.
\bibitem{HMPST} A.~Hajnal, W.~Maass, P.~Pudl\'ak, M.~Szegedy, and G.~Tur\'an.
Threshold circuits of bounded depth.
\emph{Journal of Computer and System Sciences} 46(2):129--154, 1993.
\bibitem{HP} K.~A. Hansen and V.~V. Podolskii.
Polynomial threshold functions and Boolean threshold circuits.
\emph{Information and Computation} 240:56--73, 2015.
References to theorem and lemma numbers use the
\href{https://eccc.weizmann.ac.il/report/2013/021/download/}{ECCC version, TR13-021}.
\bibitem{MP} M.~L. Minsky and S.~A. Papert.
\emph{Perceptrons: An Introduction to Computational Geometry}.
MIT Press, 1969; expanded edition, 1988. The parity-order argument is in
Section~3.1.
\bibitem{IPS} R.~Impagliazzo, R.~Paturi, and M.~Saks.
Size--depth tradeoffs for threshold circuits.
\emph{SIAM Journal on Computing} 26(3):693--707, 1997.
\bibitem{CSS} R.~Chen, R.~Santhanam, and S.~Srinivasan.
Average-Case Lower Bounds and Satisfiability Algorithms for Small Threshold
Circuits. \emph{Theory of Computing} 14(9):1--55, 2018.
\end{thebibliography}
\end{document}
